\documentclass[11pt]{article}

\usepackage[a4paper,margin=23mm,headheight=15pt]{geometry}
\usepackage{fontspec}
\usepackage{amsmath,amsthm,mathtools}
\usepackage{unicode-math}
\setmainfont{Times New Roman}
\setsansfont{Arial}
\setmonofont{Menlo}
\setmathfont{STIX Two Math}
\newfontfamily\japanesefont{Hiragino Mincho ProN}
\XeTeXlinebreaklocale "ja"
\XeTeXlinebreakskip=0pt plus 1pt
\usepackage{microtype}
\usepackage{booktabs,tabularx,longtable,array}
\usepackage{enumitem}
\usepackage{xcolor}
\usepackage{tikz}
\usetikzlibrary{arrows.meta,positioning,fit,calc}
\usepackage{fancyhdr}
\usepackage{hyperref}

\definecolor{navy}{HTML}{17365D}
\definecolor{blue}{HTML}{2F5597}
\definecolor{lightblue}{HTML}{EAF1F8}
\definecolor{warm}{HTML}{F6EEE7}
\definecolor{deepred}{HTML}{8B1E2D}
\definecolor{graytext}{HTML}{4A4A4A}
\definecolor{green}{HTML}{215E21}
\hypersetup{
  colorlinks=true,
  linkcolor=navy,
  citecolor=deepred,
  urlcolor=blue,
  pdfauthor={Satoru Watanabe},
  pdftitle={Subjectivity-Intersection Braid Quantum Gravity: Concise Research Note},
  pdfsubject={A constructive refutation of objectivist primacy from a source-derived Lorentz carrier and a conditional Einstein limit},
  pdfkeywords={subjectivity intersection, braid quantum gravity, objectivist primacy, Lorentz carrier, Einstein equation, emergent spacetime}
}

\newtheorem{theorem}{Theorem}[section]
\newtheorem{proposition}[theorem]{Proposition}
\newtheorem{lemma}[theorem]{Lemma}
\newtheorem{corollary}[theorem]{Corollary}
\theoremstyle{definition}
\newtheorem{definition}[theorem]{Definition}
\newtheorem{assumption}[theorem]{Assumption}
\newtheorem{hypothesis}[theorem]{Hypothesis}
\theoremstyle{remark}
\newtheorem{remark}[theorem]{Remark}

\newcommand{\cA}{\mathcal{A}}
\newcommand{\cE}{\mathcal{E}}
\newcommand{\cS}{\mathcal{S}}
\newcommand{\Sym}{\operatorname{Sym}}
\newcommand{\rank}{\operatorname{rank}}
\newcommand{\Tr}{\operatorname{Tr}}
\newcommand{\Ad}{\operatorname{Ad}}
\newcommand{\id}{\mathbf{1}}
\newcommand{\OPT}{\mathsf{OP}}
\newcommand{\MMR}{\mathsf{MMR}}
\newcommand{\CGR}{\mathsf{CGR}}

\pagestyle{fancy}
\fancyhf{}
\fancyhead[L]{\small\color{graytext}Braid Quantum Gravity: Concise Research Note}
\fancyhead[R]{\small\color{graytext}Accompanies full v0.7}
\fancyfoot[C]{\thepage}

\setlength{\parindent}{0pt}
\setlength{\parskip}{0.60em}
\setlength{\emergencystretch}{1.5em}
\setlist{nosep,leftmargin=1.55em}
\renewcommand{\arraystretch}{1.15}

\title{\vspace{-1.2cm}\textbf{Subjectivity-Intersection Braid Quantum Gravity}\\[0.38em]
\Large A Constructive Refutation of Objectivist Primacy\\[-0.02em]
from a Source-Derived Lorentz Carrier and a Conditional Einstein Limit\\[0.3em]
\large Concise Research Note Accompanying the Full v0.7 Preprint}
\author{Satoru Watanabe\\\small Subjectivity Intersection Emergence Lab (SIEL)}
\date{Concise research note --- 22 September 2026}

\begin{document}
\maketitle

\begin{center}
\fcolorbox{navy}{lightblue}{\parbox{0.93\linewidth}{\small
\textbf{Document and claim status.} This concise note records the post-v0.6.1 result and does not replace the full v0.7 derivation, which retains the complete v0.6.1 finite geometry, action, locality, and audit chain. The theory is a mathematically explicit \emph{candidate architecture} for quantum gravity, not a completed or empirically confirmed theory. In the declared fixed source model, the pointed-Braid event law now supplies: an operational four-direction source cylinder; a full-rank local-label matter sector; fixed-model matter routing; the relative factor $3/5$ without manual insertion; a source-null tetrad; and a Lorentz-signature event carrier whose normalization is unique in a stated naturality class. The remaining non-Braid inputs are actual source-to-smooth convergence and the \emph{Source-Dynamical Parent Completion} (SDPC). Consequently $G_{\mu\nu}=(3/5)T_{\mu\nu}$ is a conditional continuum theorem, not yet an unconditional Braid-only field equation.}}
\end{center}

\begin{abstract}
We develop Subjectivity-Intersection Braid Quantum Gravity as a candidate theory in which public spacetime structure is generated by the intersection of nonidentical, standpoint-bearing Braid processes rather than assumed as an ontologically prior container. Subjectivity here means a located mode of access, transport, or presentation, not consciousness. The theory begins with a pointed quantum-algebraic Braid source and preserves the distinction between source structure, derived mathematical structure, physical interpretation, and empirical claim.

The new result is a source-derived four-dimensional Lorentz carrier. The actual involutive Braid event transition produces a permutation-covariant event moment
\[
 Q_4=\frac{104I_4-20\mathbf1\mathbf1^{\mathsf T}}{50}
     =\frac{12}{25}P_0+\frac{52}{25}P_1,
\]
which is positive, rank four, and fixed without axis weights. Independently, the source-null causal structure selects the grading
\[
 J_{\rm ray}=I_4-\frac12\mathbf1\mathbf1^{\mathsf T}=-P_0+P_1.
\]
Their product
\[
 K_{\rm evt}=J_{\rm ray}Q_4
 =-\frac{12}{25}P_0+\frac{52}{25}P_1
\]
has exact signature $(3+,1-)$. A positive self-adjoint $S_4$-equivariant normalizer is uniquely fixed by the source-null inverse Gram tensor. No Wick rotation, manual sign choice, or continuous coefficient fit enters this carrier theorem. In parallel, the fixed event-conditioned matter model selects the full source corner, preserves ten independent response directions, closes the previously isolated matter-routing ambiguity within its declared partition class, and yields the relative source-cup normalization $3/5$ without setting that number by hand.

The result does not yet identify a complete Braid-derived action or prove convergence of the actual finite source sequence to a smooth field. The certified relative-entropy edge action remains positive and does not by itself insert $J_{\rm ray}$. Forward/reverse Braid orientation gives the same quadratic moment, while a doubled closed-time-path branch form lives on branch $\otimes$ base space and contracts only to scalar multiples of $Q_4$. These exact no-go results prevent a carrier identity from being silently promoted into a variational derivation. If one adds actual source-to-smooth convergence and the remaining Source-Dynamical Parent Completion (SDPC)---one parent action supplying action-level causal grading, metric/coframe soldering, nontrivial physical-section dynamics, and the associated Ward identity---the long-wavelength metric-affine action yields
\[
 G_{\mu\nu}(g)=\frac35T_{\mu\nu}
\]
on a continuous spacetime neighborhood. Thus the Einstein limit is conditional, whereas the finite event dynamics, matter routing, factor $3/5$, null causal geometry, and Lorentz carrier are source-derived within the fixed model.

This strengthens a constructive refutation of \emph{objectivist primacy}. A stable Lorentzian public carrier can be generated from maintained standpoint difference and intersection; it need not be posited as primitive in every admissible construction. The result does not refute mind-independent reality, establish phenomenal consciousness, or show that nature realizes the model. It provides a sharper quantum-gravity candidate and a smaller, explicitly falsifiable completion problem.
\end{abstract}

\noindent\textbf{Keywords:} Subjectivity-Intersection Braid Quantum Gravity; quantum gravity candidate; pointed braid; generated objectivity; objectivist primacy; Lorentz signature; emergent spacetime; Einstein equation; metric-affine gravity; Ward identity.

{\japanesefont
\section*{\japanesefont 日本語要旨}

本稿は、主観交差Braidを根底とする量子重力候補を提示する。ここで「主観」は意識ではなく、アクセス・輸送・表示の位置づけられた様態を意味する。異なる主観的立場を同一化せず、その交差から共有可能な時空構造が生成される、という順序を数理化する。

v0.7の中心成果は、固定source模型の内部で、Braid event則とsource因果構造からLorentz符号をもつ4方向carrierを導いたことである。event moment $Q_4$は正定値rank 4であり、sourceが選ぶ時間方向反転$J_{\rm ray}$との積$K_{\rm evt}=J_{\rm ray}Q_4$はexactに$(3+,1-)$となる。source-null Gramへの正規化も、正・自己共役・$S_4$同変という自然性クラス内で一意である。Wick回転、手動符号、連続係数fitは使わない。さらに固定event-conditioned matter模型ではMMR2 routingが閉じ、相対係数$3/5$も手動設定ではなくsource cupから得られる。

ただし、Lorentz carrierが得られたことと、実際の変分作用がそのcarrierを選ぶことは別である。現行の正のrelative-entropy作用は$J_{\rm ray}$挿入をまだ導かず、物質operator応答を保存する物理stressへ一意にsolderする則とWard恒等式も残る。したがって$G_{\mu\nu}=(3/5)T_{\mu\nu}$は現時点では条件付き連続極限定理であり、無条件Braid-only Einstein方程式ではない。本稿の主張は、客観性そのものの否定ではなく、客観的Lorentz構造を存在論的に最初から置くことが普遍的に必要だという「客観主義的優先性」への構成的反証である。
}

\section{What changed from v0.4 through v0.61}

Version 0.4 already separated an exact finite Braid construction from a conditional continuum Einstein equation. The intervening development through v0.61 refined the matter and continuum dependencies but did not yet contain SDPC. After v0.61, BGCE245 introduced \emph{Source-Dynamical Parent Completion} (SDPC) as one explicitly compound parent-action axiom replacing separate stress-solder and Ward assumptions. Subsequent certificates then tested and derived individual ingredients of SDPC without claiming the whole axiom had been derived. Version 0.7 records that post-v0.61 development and resolves the earlier bridge block into independently testable statements.

\begin{table}[ht]
\centering
\small
\begin{tabularx}{\linewidth}{>{\raggedright\arraybackslash}p{0.26\linewidth} >{\raggedright\arraybackslash}p{0.31\linewidth} X}
\toprule
Question & v0.7 result & Remaining boundary \\
\midrule
Operational four-direction carrier & Derived from the source spectral event cylinder and four null links & No claim that every possible model must have four dimensions \\
Matter realization (MMR2) & Closed in the fixed event-conditioned model and complete five-partition retraction class & Arbitrary instruments, arbitrary subalgebras, and natural-sector selection are not classified \\
Relative factor $3/5$ & Derived from normalized source-cup routing; not manually set & Its empirical value and physical-unit map are open \\
Lorentz signature & Derived algebraically from $J_{\rm ray}Q_4$, with unique natural normalizer & The certified positive edge action does not yet select the grading \\
Physical stress and conservation & Ten internal response directions exist & Unique metric/coframe solder and Ward identity remain open \\
Einstein equation & Conditional neighborhood Euler--Lagrange theorem & Unconditional Braid-only derivation remains open \\
\bottomrule
\end{tabularx}
\caption{The v0.7 dependency split. ``Derived'' always means derived inside the declared fixed source model and naturality class, not established as a universal law of nature.}
\label{tab:delta}
\end{table}

This separation changes the strongest defensible claim. The coefficient $3/5$ is no longer an inserted bridge constant, and MMR2 is no longer an independent matter-routing postulate within the fixed model. SDPC, introduced after v0.61, is the remaining named completion condition. The paper does not disguise its still-open action selection, stress solder, and Ward content as consequences of a matching signature.

\section{Subjectivity intersection and the target of the refutation}

\subsection{Subjectivity is not consciousness}

A \emph{subjectivity} in this work is a typed and located mode of access, transport, or presentation whose distinction from other such modes is constitutive of the construction. It may be represented by a source prefix, a reference system, or an oriented transport context. No braid, photon, clock, algebra, or source sector is thereby attributed sensation, awareness, intention, or phenomenal experience. Any bridge from this formal notion to consciousness would require a separate axiom and separate evidence.

Let $A$ and $B$ denote nonidentical source standpoints. Their intersection does not mean a third external observer that views both from nowhere. It is the compositional event in which their difference is retained while presentation redundancy is removed. The public structure $C(A,B)$ is then an output of the intersection rather than an independently completed container.

\begin{definition}[Generated objectivity]
A carrier $C(A,B)$ exhibits generated objectivity when it is invariant under allowed internal presentation changes, covariant under ambient re-description, symmetric under exchange of the two typed standpoint roles when the source law requires it, altered by the deletion of either role, and regenerated by one source law through the declared refinement stages without stage-dependent fitting.
\end{definition}

\subsection{The exact objectivist claim being denied}

\begin{definition}[Objectivist Primacy Thesis, $\OPT$]
Every mathematically admissible construction of a stable, shareable, transformation-respecting physical carrier must take that carrier, or an equivalent perspective-neutral carrier, as primitive. Distinct standpoints may describe it, but cannot generate it.
\end{definition}

$\OPT$ is a universal necessity thesis. It is stronger than scientific reproducibility, realism about measured regularities, covariance, or mind-independent existence. A single valid generated-objectivity model refutes $\OPT$ as stated. It does not prove that the model is realized in nature. The present paper therefore keeps two conclusions separate:
\begin{enumerate}
\item a mathematical counterexample to the necessity of objectivist primacy; and
\item a physical quantum-gravity hypothesis whose empirical truth remains open.
\end{enumerate}

The proposal is adjacent to, but not identical with, relational quantum mechanics, quantum-reference-frame programmes, QBism, perspectival realism, and structural realism \cite{rovelli1996,giacomini2019,vanrietvelde2020,fuchs2013,massimi2022,ladyman1998,weyl1952}. Its distinct technical claim is generative: a shared tensor carrier is reconstructed from maintained source provenance and intersection rather than presupposed as a perspective-neutral starting object.

\section{Pointed-Braid source and operational event cylinder}

\subsection{Quantum-algebraic source}

At refinement stage $n$, take the matrix algebra
\begin{equation}
 \cA_n=M_{5^n}(\mathbb C),\qquad
 j_n(X)=X\otimes I_5.
\end{equation}
The source contains distinguished states, marked local algebras, and unitary Braid transports satisfying the typed Yang--Baxter and composition laws. This is quantum in the algebraic sense: states, complex noncommutative observables, tensor composition, and unitary transport are present. The present paper does not derive complex amplitudes, the Born rule, or Hilbert-space composition from an unstructured topological braid group.

The source data are therefore neither ``bare braid topology'' \cite{artin1947} nor a continuum spacetime. They are a finite quantum-algebraic pointed-Braid source developed in the preceding Subjectivity-Intersection Mathematics programme \cite{watanabe2026braid,watanabe2026ontology}. The scientific question is how much geometry and dynamics this source forces without target fitting.

\subsection{Operational event identification}

The fixed source packet contains a marked event effect $P$ and a spectral event algebra. The actual stored event transition $R_{\rm evt}$ acts on 25 local atoms, is involutive, and satisfies the declared Yang--Baxter event relation. The same marked generator determines four future-null affine links. This identification retires the independent \emph{finite event-algebra correspondence} inside the fixed model: the event carrier is read from the source law itself rather than attached as a separate map. It does not retire the source-to-continuum part of CGR.

This is a scoped statement. It does not say that every operational event in nature must be represented by this $P$, nor does it select a natural signed sector. It says only that the previously separate finite event correspondence is redundant in the declared source construction. Convergence of the actual finite Braid coframe/connection sequence to a uniformly regular smooth field, including convergence of first variations, remains a separate assumption.

\subsection{Four dimensions without a primitive completed spacetime}

The operational cylinder has four source directions because four independently certified null links span the event carrier. Let $L$ be the matrix with these links as columns. The exact certificate gives
\begin{equation}
 \det L=\frac{81}{16}\neq0.
\end{equation}
Thus the event carrier is four-dimensional in the fixed model. This is stronger than postulating an arbitrary ambient $\mathbb R^4$ and then placing Braid data inside it: the four operational directions are generated and are nondegenerate. It is weaker than a theorem that four dimensions are inevitable for all admissible Braid sources. The latter is not claimed.

\section{Matter routing and the non-manual factor $3/5$}

\subsection{The full event-conditioned corner}

The marked event defines the physical corner candidate
\begin{equation}
 \cA_{\cap}=P\cA P,
\end{equation}
with Hamiltonian $H_{\cap}=PG_1P$ and ten compressed source operators
\begin{equation}
 F_A^{\cap}=PF_AP,\qquad A=1,\ldots,10.
\end{equation}
The corresponding Gibbs family is faithful on the corner, and its Bogoliubov--Kubo--Mori Hessian has rank ten. Hence the ten declared matter-response directions are locally nondegenerate.

The source atoms generate five partition retractions. Requiring exact preservation of the already fixed eleven-operator packet
\begin{equation}
 \cS_{\cap}=\{PG_1P,PF_1P,\ldots,PF_{10}P\}
\end{equation}
selects only the full partition. Every nontrivial dephasing changes $H_{\cap}$ before any coefficient can be retuned. This closes $\MMR2$ within the declared fixed event-conditioned model and complete five-partition class. It does not prove uniqueness among all completely positive instruments or all possible subalgebras.

\subsection{Where $3/5$ comes from}

The normalized tail trace sends the corner unit to the outer event effect $P$. Its source-cup composition yields the relative multiplier
\begin{equation}
 \kappa_B=\frac35.
\end{equation}
No line of the v0.7 Lorentz-carrier calculation sets or retunes this value. It is inherited from the same marked source routing that selects the event-conditioned matter packet. In particular, agreement with the coefficient appearing in the conditional Einstein equation is not counted as proof of that equation: a source-derived normalization is still distinct from a conserved physical stress tensor.

\section{Braid-derived local-label dynamics}

\subsection{Actual event transition versus identity and flip}

For each adjacent pair of source directions, define the displacement under the actual Braid event transition. Its exact pair moment is
\begin{equation}
 M_{\rm pair}=
 \begin{pmatrix}
 28&-20\\[-0.15em]
 -20&28
 \end{pmatrix},
\end{equation}
with eigenvalues $8$ and $48$ and rank two. The identity comparator has rank zero and the ordinary flip comparator rank one. Thus the rank-two coupling is specific to the actual non-flip event map within this comparator class.

Average the adjacent-pair moment over the 24 source charts. The result is
\begin{equation}\label{eq:q4}
 Q_4=\frac{104I_4-20\mathbf1\mathbf1^{\mathsf T}}{50}
 =\frac{12}{25}P_0+\frac{52}{25}P_1,
\end{equation}
where
\begin{equation}
 P_0=\frac14\mathbf1\mathbf1^{\mathsf T},\qquad
 P_1=I_4-P_0.
\end{equation}
$Q_4$ has diagonal entries $42/25$, off-diagonal entries $-2/5$, and eigenvalues $12/25$ on the uniform line and $52/25$ on its three-dimensional complement. It is positive and full rank. No distinguished physical $1+3$ ordering or external axis weights enter Eq.~\eqref{eq:q4}.

\subsection{Relative entropy and the BKM--Dirichlet form}

The source response is the derivative of a log-partition function. Within the response-exact Bregman class this selects the log-$Z$ Bregman divergence, equivalently the Umegaki relative entropy for the Gibbs family \cite{umegaki1962,amari2016}. On actual neighbor event edges, the quadratic term is
\begin{equation}\label{eq:bkm}
 \delta^2 S_{\rm edge}
 =\frac12\int Q_4^{ab},
 \langle \partial_a\lambda,\partial_b\lambda\rangle_{\rm BKM},d\mu.
\end{equation}
The base moment has rank four and the internal BKM metric rank four in the tested local-label sector, so the gradient quadratic form has rank sixteen. This supplies source-independent local-label dynamics on the fixed cylinder. It is a positive Dirichlet form, not yet a Lorentzian action.

\section{Null tetrad and source-derived Lorentz carrier}

\subsection{The source-null metric}

The four null links have a source-fixed Gram matrix $G_{\rm null}$ with zero diagonal and off-diagonal entries $-12/25$. Since $L$ is invertible, the ambient metric reconstructed from the links is unique:
\begin{equation}\label{eq:g-null}
 g=L^{-\mathsf T}G_{\rm null}L^{-1}
 =\operatorname{diag}\!\left(
 -\frac{16}{25},\frac{16}{75},\frac{16}{75},\frac{16}{75}
 \right).
\end{equation}
It has signature $(3+,1-)$. The induced map on symmetric tensors, $\Sym^2(L)$, has rank ten. Thus the four-direction carrier is algebraically capable of supporting all ten symmetric rank-two components.

This statement alone does not solve the physical stress problem. Full rank is an existence fact about an algebraic map; it does not prove that the actual matter operator packet transforms through that map or that the resulting response is conserved.

\subsection{Why a coframe alone cannot create the Lorentz sign}

The positive event moment remains positive under the null coframe:
\begin{equation}
 LQ_4L^{\mathsf T}
 =\operatorname{diag}\!\left(
 \frac{27}{25},\frac{117}{25},\frac{117}{25},\frac{117}{25}
 \right).
\end{equation}
By Sylvester's law of inertia, no real invertible congruence changes signature $(4+,0-)$ to $(3+,1-)$. Therefore a claim that ``the null tetrad automatically Lorentzianizes the positive edge action'' is false. This no-go forces the causal sign to have a source origin other than coordinate change.

\subsection{Causal grading}

The source causal analysis distinguishes the uniform sum of the four future-null rays as the timelike direction. The corresponding reflection is
\begin{equation}\label{eq:jray}
 J_{\rm ray}=I_4-\frac12\mathbf1\mathbf1^{\mathsf T}
 =-P_0+P_1.
\end{equation}
It has eigenvalue $-1$ on the uniform line and $+1$ on the standard three-dimensional complement. Since $Q_4$ is $S_4$-equivariant, $J_{\rm ray}$ and $Q_4$ commute. Hence
\begin{equation}\label{eq:kevt}
 K_{\rm evt}=J_{\rm ray}Q_4
 =-\frac{12}{25}P_0+\frac{52}{25}P_1
\end{equation}
is symmetric with exact signature $(3+,1-)$. In components its diagonal is $36/25$ and every off-diagonal entry is $-16/25$.

\begin{theorem}[Source-derived Lorentz event carrier]\label{thm:carrier}
In the declared fixed source model, the actual Braid event moment $Q_4$ and the independently source-selected causal grading $J_{\rm ray}$ determine the Lorentz-signature carrier $K_{\rm evt}=J_{\rm ray}Q_4$. The construction uses no Wick rotation, manual sign assignment, or fitted continuous coefficient.
\end{theorem}

\subsection{Natural normalization to the null Gram tensor}

Because Eq.~\eqref{eq:kevt} is contravariant in the kinetic form, its physical comparison target is $G_{\rm null}^{-1}$, not $G_{\rm null}$. Restrict a normalizer to the positive, self-adjoint, $S_4$-equivariant class
\begin{equation}
 A_{\rm kin}=a_0P_0+a_1P_1.
\end{equation}
Then
\begin{equation}
 A_{\rm kin}^{\mathsf T}K_{\rm evt}A_{\rm kin}=G_{\rm null}^{-1}
\end{equation}
fixes
\begin{equation}
 a_0=\frac{25\sqrt3}{36},\qquad
 a_1=\frac{25\sqrt{39}}{156}
\end{equation}
uniquely within that class. All 24 chart permutations commute with the projector decomposition. In the character frame one recovers
\begin{equation}
 g^{-1}=\operatorname{diag}\!\left(
 -\frac{25}{16},\frac{75}{16},\frac{75}{16},\frac{75}{16}
 \right).
\end{equation}
The uniqueness claim is deliberately local: it is not uniqueness under arbitrary Lorentz isometries or among all nonlinear field redefinitions.

\section{What the carrier theorem does not yet prove}

\subsection{Carrier selection is not action selection}

The certified edge action underlying Eq.~\eqref{eq:bkm} is a positive sum of Umegaki relative entropies. It yields $Q_4$, not $J_{\rm ray}Q_4$. Both factors in $K_{\rm evt}$ are source-derived, but it remains a separate proposition that the physical variational principle multiplies them in precisely this way. Inserting Eq.~\eqref{eq:jray} by declaration would produce a Lorentzian action, but would be a new bridge axiom rather than a derivation.

\subsection{Two exact failed shortcuts}

The first shortcut uses Braid orientation. Forward and reverse event displacements have the same second moment. Even and odd chart permutations also average to the same $Q_4$. Their signed difference is zero, and every scalar linear combination of forward and reverse moments is only a scalar multiple of $Q_4$. Orientation therefore survives in the word data but is erased by the quadratic relative-entropy Hessian.

The second shortcut uses doubled closed-time-path branches. The branch form
\begin{equation}
 \eta_{\rm br}\otimes Q_4,
 \qquad \eta_{\rm br}=\operatorname{diag}(1,-1),
\end{equation}
has signature $(4+,4-)$ on branch $\otimes$ base space. A fixed branch contraction gives
\begin{equation}
 (v^{\mathsf T}\eta_{\rm br}v)Q_4,
\end{equation}
which is $Q_4$, $-Q_4$, zero, or another scalar multiple. It cannot yield the mode-selective $K_{\rm evt}$. A branch-to-ray intertwiner designed to do so would restate the missing assumption.

These no-go results matter scientifically. They prevent a visually attractive $(3+1)$ signature from being counted twice: once as a carrier theorem and again, without proof, as an action theorem.

\section{Conditional continuum Einstein limit}

\subsection{The remaining bridge: Source-Dynamical Parent Completion}

Earlier versions listed four continuum laws: temporal soldering, continuum/infrared convergence, matter response with Ward identity, and common normalization. The v0.7 dependency audit changes that ledger. The operational four-direction carrier, fixed-model matter routing, and common factor $3/5$ now have source derivations in the declared model. Two gates remain distinct: (i) actual source-to-smooth convergence with first-variation control, and (ii) the following compound but nontrivial parent-action principle.

\begin{assumption}[Source-Dynamical Parent Completion, SDPC]\label{ass:sdpc}
On the same source cylinder there exists one local parent matter action $S_{\rm parent}[g,\Psi;a]$ such that:
\begin{enumerate}
\item its source first jet reproduces the already derived full-corner responses and fixed-model MMR2 routing;
\item its metric/coframe coupling fixes the pullback of the ten source derivatives, eliminating an externally chosen solder and its scale ambiguity;
\item it contains nontrivial source-independent dynamics on the physical section $a=0$, and its kinetic principal symbol selects $J_{\rm ray}Q_4$ rather than inserting $J_{\rm ray}$ after variation;
\item it is invariant under the diagonal diffeomorphism/frame action, so one off-shell Noether identity links stress divergence to the matter Euler--Lagrange terms;
\item the same matter fields are placed on shell by this action, yielding covariant stress conservation; and
\item in the low-curvature, matter-hypermomentum-free sector, the compatible Palatini connection equation selects the Levi--Civita connection up to removed projective freedom.
\end{enumerate}
\end{assumption}

SDPC is the compound axiom isolated in BGCE245, not a new v0.7 replacement name. It compresses the previously separate stress-solder and Ward assumptions because both must follow from the same parent action. Later certificates derive important ingredients---the source first jet, fixed-cylinder local-label dynamics, the functional form of the positive edge action, and the Lorentz carrier---but do not derive the whole conjunction. In particular, the positive edge action still fails to select the signed grading. SDPC therefore remains open as a Braid-only theorem.

\subsection{Variational derivation under SDPC}

Under SDPC, the long-wavelength metric-affine action takes the standard first-order form familiar from metric-affine gravity \cite{hehl1995},
\begin{equation}\label{eq:sir}
 S_{\rm IR}[g,\Gamma,\psi]
 =\frac{1}{2\kappa_B}\int d^4x\,\sqrt{-g}\,
 g^{\mu\nu}R_{(\mu\nu)}(\Gamma)
 +S_m[g,\Gamma,\psi],
 \qquad \kappa_B=\frac35.
\end{equation}
The stress tensor is defined by the variation of the same matter action,
\begin{equation}\label{eq:stress}
 T_{\mu\nu}=-\frac{2}{\sqrt{-g}}
 \frac{\delta S_m}{\delta g^{\mu\nu}}.
\end{equation}
Metric or coframe variation gives
\begin{equation}
 R_{(\mu\nu)}(\Gamma)-\frac12g_{\mu\nu}R(\Gamma)
 =\kappa_BT_{\mu\nu}.
\end{equation}
The SDPC-compatible connection equation then selects $\Gamma=\Gamma_{\rm LC}(g)$ in the declared infrared sector. Substitution yields
\begin{equation}\label{eq:einstein}
 \boxed{G_{\mu\nu}(g)=\frac35T_{\mu\nu}}.
\end{equation}

\begin{theorem}[Conditional continuum Einstein theorem]\label{thm:einstein}
The exact fixed-model source results summarized in Sections 3--6, together with actual source-to-smooth convergence and SDPC, imply Eq.~\eqref{eq:einstein} as an Euler--Lagrange equation on a continuous spacetime neighborhood in the long-wavelength, low-curvature, matter-hypermomentum-free sector.
\end{theorem}

The theorem is more than a pointwise metric fit: it is an off-shell variational statement on a neighborhood. It is less than an unconditional Braid-only derivation because neither actual source-to-smooth convergence nor SDPC is yet a theorem of the source.

\subsection{The exact status of $3/5$}

The number $3/5$ in Eq.~\eqref{eq:einstein} is source-derived within the fixed event-conditioned matter model. The conditionality of the equation no longer comes from choosing this number. It comes from identifying the source response with the conserved physical stress of the same variational theory. This distinction is essential: coefficient derivation does not by itself establish tensor typing, conservation, or dynamics.

\section{Is this already quantum gravity?}

The strongest accurate description is \emph{a source-explicit quantum-gravity candidate with an exact finite Lorentz carrier and a conditional Einstein limit}. That phrase has four parts.

\begin{enumerate}
\item \textbf{Quantum source.} The microscopic source is a finite complex noncommutative algebra with states, unitary Braid transport, and tensor refinement.
\item \textbf{Geometry from the source.} Four operational directions, null geometry, a Lorentz carrier, finite coframe/metric-affine updates, and nonzero curvature are outputs in the declared model.
\item \textbf{Matter from the same event structure.} The full corner action, ten response directions, MMR2 routing, and relative factor $3/5$ are selected without target-Einstein fitting inside the fixed class.
\item \textbf{Gravity limit.} An Einstein equation follows variationally under actual source-to-smooth convergence and SDPC, but both the continuum limit and source-only action/stress/Ward closure are unfinished.
\end{enumerate}

This exceeds a mere analogy between braids and spacetime. It does not yet meet the standard of a completed fundamental theory. In particular, the paper does not provide a continuum graviton Hilbert space, renormalization or ultraviolet completion, anomaly-free first-class constraint algebra, black-hole or cosmological solutions uniquely predicted by the source, calibrated Newton constant, Standard Model coupling, or natural-data confirmation.

\section{Why subjectivity intersection matters}

Could an ordinary braid be substituted with no loss? Not for the present derivation as written. The construction uses more than braid-group relations:
\begin{enumerate}
\item a pointed event effect $P$ selects an operational corner;
\item nonidentical source prefixes retain provenance through transport;
\item the intersection packet is tested against collapse of either standpoint;
\item source coherence, not only algebraic rank, selects the full matter corner;
\item the common causal grading and event moment arise from the same marked source lineage.
\end{enumerate}
An unpointed braid group does not by itself supply these typed roles. One could add equivalent markings, states, and provenance to an ``ordinary'' braid construction, but then one would have recreated the operative subjectivity-intersection structure under another name.

This does not establish that the philosophical interpretation is unique. The exact mathematics can also be described in the language of marked quantum events, relational reference structures, or source-provenance-preserving transport. The specifically SIEL claim is that these formal roles realize a minimal grammar of distinct standpoints generating a public carrier. The empirical task is to determine whether this interpretation has explanatory or predictive force beyond the mathematics.

\section{Constructive refutation of objectivist primacy}

The counterexample now has the following form:
\begin{equation}
 \begin{aligned}
 &\text{nonidentical pointed source standpoints}\\
 &\quad\longrightarrow\text{operational Braid event and full matter corner}\\
 &\quad\longrightarrow\text{four source-null directions and causal grading}\\
 &\quad\longrightarrow\text{shared Lorentz carrier with fixed natural scale}.
 \end{aligned}
\end{equation}
The public Lorentz carrier is stable, chart-covariant in the declared $S_4$ class, full rank, and reproducible. Deleting the marked structure or replacing the event transition with identity/flip comparators destroys essential parts of the construction. Thus the carrier is neither a private appearance nor a primitive completed background. It is generated by a relational source with maintained provenance.

This refutes $\OPT$: it is false that every admissible stable Lorentzian public carrier must be ontologically primitive. It does not refute the existence of objective structure. On the contrary, it explains objectivity as invariance, covariance, reproducibility, and shared accessibility after intersection. Nor does it establish that our universe is generated by this mechanism. The mathematical and physical claims must remain distinct.

\section{Falsifiability and completion gates}

\subsection{Formal falsifiers}

The present exact package fails if any of the following occurs under revision-matched replay:
\begin{itemize}
\item the actual event map loses its rank-two pair moment or fails the stored Yang--Baxter relation;
\item the $S_4$ average depends on a hidden physical axis weight;
\item a nontrivial source partition preserves the fixed eleven-operator matter packet exactly;
\item the normalized source routing no longer yields $3/5$ without retuning;
\item the four null links become linearly dependent or fail to reconstruct Eq.~\eqref{eq:g-null};
\item $J_{\rm ray}$ ceases to be source-selected or $K_{\rm evt}$ loses signature $(3+,1-)$;
\item the claimed natural normalizer is nonunique in the declared positive self-adjoint $S_4$-equivariant class; or
\item exact certificate regeneration ceases to reproduce the frozen outputs.
\end{itemize}

\subsection{The shortest completion gate}

The shortest route to an unconditional field equation is not another coefficient scan. It is a source derivation of SDPC. The next decisive test should construct a first-order phase-space or Legendre action directly from the marked Braid generator and ask whether its principal symbol acts as $-P_0+P_1$ on the same event-gradient field while its source derivatives reproduce the already fixed full-corner responses. A passing result must simultaneously provide:
\begin{enumerate}
\item action-level selection of $J_{\rm ray}$;
\item a typed map from the ten source derivatives to metric/coframe variations;
\item an off-shell Ward identity; and
\item preservation of the source-derived $3/5$ without readjustment.
\end{enumerate}
Failure of any item blocks promotion to Braid-only Einstein gravity. Standardly postulating $T-V$, a symplectic form, or the Hilbert stress would not count as derivation unless the same source fixes it.

\subsection{Empirical gate}

Even a complete SDPC derivation would not by itself show that nature realizes the model. A physical theory requires units, scale matching, instrument models, noise and uncertainty, and at least one discriminator against general relativity plus ordinary quantum matter or against lower-commitment relational alternatives. No such natural-data discriminator is claimed here.

\section{Claim ledger}

\begin{longtable}{>{\raggedright\arraybackslash}p{0.35\linewidth} >{\raggedright\arraybackslash}p{0.18\linewidth} >{\raggedright\arraybackslash}p{0.38\linewidth}}
\toprule
Statement & Status & Scope \\
\midrule
\endfirsthead
\toprule
Statement & Status & Scope \\
\midrule
\endhead
Operational event carrier from the marked source & Exact derived & Fixed source model \\
Finite event-algebra correspondence & Retired & Redundant inside the fixed operational identification \\
Full $\CGR$ source-to-continuum statement & Open & Actual finite-source convergence and first-variation control are not proved \\
Full-corner matter action and rank-ten response & Exact derived & Fixed event-conditioned corner \\
$\MMR2$ matter routing & Closed & Complete five-partition retraction class; not arbitrary instruments \\
Relative factor $3/5$ & Source-derived & Fixed normalized source-cup routing \\
Braid event moment $Q_4$ & Exact derived & 24-chart $S_4$ average \\
Four-dimensional null carrier & Exact derived & Four certified links; not universal dimension uniqueness \\
Lorentz carrier $K_{\rm evt}$ & Exact derived & Source grading times event moment \\
Normalization to $G_{\rm null}^{-1}$ & Unique in class & Positive self-adjoint $S_4$-equivariant normalizers \\
Action-level selection of $J_{\rm ray}$ & Open & Orientation-only and fixed-branch shortcuts fail \\
Physical stress solder and Ward identity & Open & Required for unconditional sourced Einstein equation \\
SDPC as one source-derived parent action & Open & Several ingredients derived; action-level grading, physical pullback, and Ward closure remain \\
Einstein equation on a neighborhood & Conditional theorem & Source-to-smooth convergence, SDPC, and the long-wavelength, low-curvature, zero-hypermomentum sector \\
Completed quantum gravity & Not established & Quantized continuum gravity and UV completion absent \\
Empirical gravity & Not established & No calibrated natural-data test \\
Constructive refutation of $\OPT$ & Established in-model & Refutes universal necessity, not mind-independent reality \\
Consciousness claim & Not made & Subjectivity is not phenomenal consciousness \\
\bottomrule
\end{longtable}

\section{Conclusion}

Subjectivity-Intersection Braid Quantum Gravity now has a more sharply divided architecture. The fixed pointed-Braid source supplies an operational event cylinder, a nondegenerate event-conditioned matter sector, source routing, the factor $3/5$, four null directions, a causal grading, and a Lorentz-signature event carrier. The carrier's scale is uniquely tied to the inverse null Gram tensor within a declared naturality class. These are not obtained by fitting an Einstein target, choosing a Wick rotation, or manually inserting $3/5$.

The central unfinished step is equally precise. The positive discrete relative-entropy action has not yet been proved to select the causal grading; the source operator responses have not yet been uniquely identified with a conserved physical stress variation; and the associated Ward identity remains open. Orientation and closed-time-path doubling do not solve this problem. Under actual source-to-smooth convergence and one explicit compound variational closure principle, the continuum metric-affine action yields $G_{\mu\nu}=(3/5)T_{\mu\nu}$ on a neighborhood. Without both conditions, the result is a source-derived Lorentz carrier and matter normalization, not an unconditional Einstein equation.

The philosophical consequence is constructive rather than eliminative. Objective Lorentz structure need not be the first ontological item in every exact model. It can arise as the stable public result of nonidentical standpoint-bearing processes intersecting under a common Braid law. Objectivity survives; its claimed universal priority does not. The physical consequence is a serious but unfinished quantum-gravity candidate whose next theorem has a clear target and whose failure conditions are explicit.

\section*{Data, code, and provenance}

The exact certificates and replay scripts are stored in the SIEL Research Agent repository,
\href{https://github.com/SIEL-Research/SIEL-Research-Agent}{SIEL-Research/SIEL-Research-Agent}. The v0.7 update uses the verified lineage through BGCE260, including the operational event identification (BGCE139), fixed-model matter selection and MMR2 closure (BGCE234--BGCE239), conditional source-derivative parent action audit (BGCE245), relative-entropy/BKM action selection (BGCE252--BGCE254), event moment and null-tetrad analysis (BGCE256--BGCE257), source-derived Lorentz carrier and normalizer (BGCE258), and the orientation and branch-space no-go results (BGCE259--BGCE260). The source revision at drafting was commit \texttt{eb61e91a0}. These are internal exact theoretical and computational certificates, not independent replication, peer review, or empirical validation.

\section*{Declaration on computational assistance}

OpenAI Codex was used under the author's direction for mathematical ideation, exact-check scripting, audit organization, and drafting. The author retains responsibility for definitions, interpretation, verification, citations, and any decision to release or submit the work.

\begin{thebibliography}{99}
\raggedright
\small

\bibitem{artin1947}
E. Artin, ``Theory of braids,'' \emph{Annals of Mathematics} \textbf{48}(1), 101--126 (1947). \href{https://doi.org/10.2307/1969218}{doi:10.2307/1969218}.

\bibitem{amari2016}
S.-i. Amari, \emph{Information Geometry and Its Applications}. Springer (2016). \href{https://doi.org/10.1007/978-4-431-55978-8}{doi:10.1007/978-4-431-55978-8}.

\bibitem{fuchs2013}
C. A. Fuchs and R. Schack, ``Quantum-Bayesian coherence,'' \emph{Reviews of Modern Physics} \textbf{85}, 1693--1715 (2013). \href{https://doi.org/10.1103/RevModPhys.85.1693}{doi:10.1103/RevModPhys.85.1693}.

\bibitem{giacomini2019}
F. Giacomini, E. Castro-Ruiz, and \v{C}. Brukner, ``Quantum mechanics and the covariance of physical laws in quantum reference frames,'' \emph{Nature Communications} \textbf{10}, 494 (2019). \href{https://doi.org/10.1038/s41467-018-08155-0}{doi:10.1038/s41467-018-08155-0}.

\bibitem{hehl1995}
F. W. Hehl, J. D. McCrea, E. W. Mielke, and Y. Ne'eman, ``Metric-affine gauge theory of gravity,'' \emph{Physics Reports} \textbf{258}, 1--171 (1995). \href{https://doi.org/10.1016/0370-1573(94)00111-F}{doi:10.1016/0370-1573(94)00111-F}.

\bibitem{ladyman1998}
J. Ladyman, ``What is structural realism?'' \emph{Studies in History and Philosophy of Science Part A} \textbf{29}(3), 409--424 (1998). \href{https://doi.org/10.1016/S0039-3681(98)80129-5}{doi:10.1016/S0039-3681(98)80129-5}.

\bibitem{massimi2022}
M. Massimi, \emph{Perspectival Realism}. Oxford University Press (2022). \href{https://doi.org/10.1093/oso/9780197555620.001.0001}{doi:10.1093/oso/9780197555620.001.0001}.

\bibitem{rovelli1996}
C. Rovelli, ``Relational quantum mechanics,'' \emph{International Journal of Theoretical Physics} \textbf{35}, 1637--1678 (1996). \href{https://doi.org/10.1007/BF02302261}{doi:10.1007/BF02302261}.

\bibitem{umegaki1962}
H. Umegaki, ``Conditional expectation in an operator algebra, IV,'' \emph{Kodai Mathematical Seminar Reports} \textbf{14}, 59--85 (1962). \href{https://doi.org/10.2996/kmj/1138844604}{doi:10.2996/kmj/1138844604}.

\bibitem{vanrietvelde2020}
A. Vanrietvelde, P. A. H\"ohn, F. Giacomini, and E. Castro-Ruiz, ``A change of perspective: switching quantum reference frames via a perspective-neutral framework,'' \emph{Quantum} \textbf{4}, 225 (2020). \href{https://doi.org/10.22331/q-2020-01-27-225}{doi:10.22331/q-2020-01-27-225}.

\bibitem{weyl1952}
H. Weyl, \emph{Symmetry}. Princeton University Press (1952).

\bibitem{watanabe2026braid}
S. Watanabe, \emph{Subjectivity-Intersection Mathematics: All-Orders Intersection Consistency and Compositional Objectivity from Pointed Unitary Braids}, version 1.0, Zenodo (2026). \href{https://doi.org/10.5281/zenodo.22166167}{doi:10.5281/zenodo.22166167}.

\bibitem{watanabe2026ontology}
S. Watanabe, \emph{Subjectivity Intersection Ontology: A Comparative Ontological Framework for Third Subjectivity}, version 1.0.0, Zenodo (2026). \href{https://doi.org/10.5281/zenodo.21641878}{doi:10.5281/zenodo.21641878}.

\end{thebibliography}

\end{document}
